\documentclass{article}
\usepackage{graphicx} % Required for inserting images
\usepackage{datetime} % Add this line to include the datetime package
\usepackage{amsmath}
\usepackage{amssymb}
\usepackage{amsthm}
\usepackage{tikz}
%\usepackage{hyperref}
\usepackage{xcolor} % Required for defining custom colors

\definecolor{darkblue}{rgb}{0.0, 0.0, 0.5} % Define a darker blue color

\usepackage[colorlinks=true, linkcolor=darkblue, urlcolor=darkblue, citecolor=darkblue]{hyperref} % Use the darker blue color

% Define theorem style
\theoremstyle{plain} % Bold title, italic body
\newtheorem{theorem}{Theorem}
\newtheorem{lemma}[theorem]{Lemma}
\newtheorem{corollary}[theorem]{Corollary}

\theoremstyle{definition} % Bold title, normal body
\newtheorem{definition}[theorem]{Definition}
\newtheorem{example}[theorem]{Example}


\title{Homework 5}
\author{Aaron Honculada}
\date{\monthname[\month] \the\year}

\begin{document}

\maketitle

\section*{Problem 1}
\noindent\textit{Recall $H = \{e, (1\text{ }2)(3\text{ }4), (1\text{ }3)(2\text{ }4),
(1\text{ }4)(2\text{ }3)\} \trianglelefteq S_4$, and $S_4/H \cong S_3$}.

\noindent(a) Find all subgroups of $S_3$ and their orders.

Both the trivial group $\{ e \}$ and $S_3$ are subgroups of $S_3$ with orders
1 and 6 respectively. There are four non-trivial proper subgroups of $S_3$. 
\begin{itemize}
    \item $Z_2 \cong \{e, (1\text{ }2)\}$ with order 2
    \item $Z_2 \cong \{e, (1\text{ }3)\}$ with order 2
    \item $Z_2 \cong \{e, (2\text{ }3)\}$ with order 2
    \item $A_3 \cong \{e, (1\text{ }2\text{ }3), (1\text{ }3\text{ }2) \}$ with order 3
\end{itemize}

\noindent(b) Using part (a), determine how many subgroups of $S_4$ contain $H$ and find their orders.

By the fourth Isomorphism Theorem, we know that for every subgroup in
$S_4/H \cong S_3$ there is a one-to-one correspondence with a subgroup of $S_4$ that contains
$H$. From this one-to-one correspondence, the number of subgroups
is still six. The three subgroups in $S_4$ that correspond to the three
subgroups in $S_3$ that are isomorphic to $\mathbb{Z}_2$ have order 8. The one
subgroup $A_4 \subset S_4$ which corresponds to $A_3 \subset S_3$ has order 12. The subgroup
corresponding to the identity in $S_3$ is the Klein-four group, which has order 4. 
The subgroup of all $S_3$ corresponds to all of $S_4$ which has order 24.

\noindent(c) Which of the subgroups from part (b) are normal?

The subgroups Klein-four group and $A_4$ are normal in $S_4$.

\section*{Problem 2}
\noindent\textit{Let $H$ and $N$ be subgroups of $G$ with $N$ normal.}

\noindent(a) Show that $N \cap H$ is a normal subgroup of $H$, and $N$ is a normal
subgroup of $NH = \{nh | n \in N, h \in H \}$.

\begin{proof}
    Let $h \in H$ and $x \in H \cap N$. Since $H$ is a subgroup of
    $G$ then we also know that $h \in G$. Given that $N$ is normal in
    $G$ then for all $n \in N$ we know that $h^{-1}nh \in N$. By definition
    of subgroup, we know that $H$ is closed and thus $h^{-1}xh \in H$.
    That means $h^{-1}xh$ is also in $N$ by normality and thus
    $h^{-1}xh \in H \cap N$.
    Since $H$ is a subgroup of $G$, then $h^{-1}xh \in H \cap N$ and 
    $H \cap N$ is normal in $H$.

    We know that $N$ is normal in $G$ so for $n \in N$ and $h \in H$ we have
    $h^{-1}nh \in N$. Therefore,
    \begin{center}
    $hN = h(h^{-1}nh) = (hh^{-1})nh = nh \in Nh$
    \end{center}
    Hence, $hN = Nh$
\end{proof}

\noindent(b) Show that $H/(N \cap H) \cong NH/N$.

\begin{proof}
    Let $\varphi: H \to HN/N$ where $\varphi(h) = hN$. We see that
    the coset $hN$ is an element of $HN/N$. Next we show that
    $\varphi$ is a surjective homomorphism.
    \begin{center}
        $\varphi(xy) = xyN = (xN)(yN) = \varphi(x)\varphi(y)$
    \end{center}
    Therefore $\varphi$ is a homomorphism. Also 
    \begin{center}
        $\varphi(h) = hN = hnN \in HN/N$
    \end{center}
    We have shown that $\varphi$ is a surjective isomorphism.
    We also observe that ker$(\varphi) = \{h \in H | \varphi(h) = e_{HN/N}\}$
    where $h$ is both an element of $H$ and of the coset $hN$. Since we have that
    $h \in H$ and $h \in N$ then ker$\varphi = H \cap N$.

    We conclude by the First Isomorphism Theorem, $H/(N \cap H) 
    \cong NH/N$.
\end{proof}

\section*{Problem 3}
\noindent\textit{Given $g,h \in G$, the commutator of $g$ and $h$ is
defined to be $[g,h] = g^{-1}h^{-1}gh$. Let $H$ be the subgroup of
$G$ generated by all commutators $[g,h]$ for $g,h\in G$.}

\noindent(a) Prove that $H$ is a normal subgroup of $G$.

\begin{proof}
It is given that $H$ is a subgroup of $G$. We now show normality of $H$
in $G$. Let $g^{-1}h^{-1}gh \in H$. We want to show that $gHg^{-1} = H$ for all $g \in G$.
\begin{center}
    $gHg^{-1} = g(g^{-1}h^{-1}gh)g^{-1} = h^{-1}ghg^{-1} \in H$
\end{center}

We have thus shown that $H$ is normal in $G$.
\end{proof}

\noindent(b) Prove that $G/H$ is abelian.

\begin{proof}
    We have that $ghg^{-1}h^{-1} \in H$ and we have shown that $H$ is normal
    in $G$. We now show that if $H$ is normal
    in $G$ then $G/H$ is abelian if and only if $ghg^{-1}g^{-1} \in H$ for all
    $g, h \in G$. We know that $G/H$ is abelian if and only if $Hgh = HgHh = HhHg = Hhg$ for
    all $g,h \in G$. We observe that $Hgh = Hhg$ if and only if $(gh)(hg)^{-1} \in H$ and 
    $(gh)(hg)^{-1} = ghg^{-1}h^{-1}$. Then it follows that since 
    $ghg^{-1}h^{-1} \in H$ for all $g,h \in G$, then $G/H$ must be abelian.
\end{proof}

\noindent(c) Show that every subgroup $N$ of $G$ containing $H$ is normal.

\begin{proof}
    To show normality we must show that $gN = Ng$ for all $g \in G$.
    Let $gn \in gN$.

    \begin{align*}
        gn &= gn(e) \\
        &= gn(g^{-1}n^{-1}ng) \\
        &= (gng^{-1}n^{-1})ng \\
        &= n'ng \\
        &= n''g \in Ng
    \end{align*}
    Therefore $gN = Ng$ and $N \supseteq H$ is normal in $G$. 
\end{proof}

\section*{Problem 4}
\noindent\textit{For $n \geq 3$, is $S_n$ isomorphic to
$A_n \times Z_2$?}
\begin{proof}
    We show that for all $n \geq 3, A_n \times Z_2 \not \cong S_n$.
    We know that $S_n$ has a trivial center, meaning that no elements in $S_n$
    other than the identity commute with every element of $S_n$. However when considering
    the element $(e, 1) \in A_n \times \mathbb{Z}_2$ we see that it is order 2, meaning that
    it is not the identity,and it commutes with every element in $A_n \times \mathbb{Z}_2$, showing that it does not
    have a trivial center. Therefore we conclude that $S_n$ is not isomorphic to $A_n \times \mathbb{Z}_2$ for $n \geq 3$.
\end{proof}

\section*{Problem 5}
\noindent\textit{Let $D_6$ be the dihedral group obtained as the symmetry
group of a regular hexagon centered at the origin with one vertex at (1, 0).
Let $r$ be the rotation by $\frac{\pi}{3}$ counterclockwise and $s$ the reflection 
over the x-axis, the elements of $D_6$ can be written as $r^is^j$, where
$0 \leq i \leq 5$ and $0 \leq s \leq 1$.}

\noindent(a) Express $D_6$ as the (internal) direct product of two nontrivial
subgroups $M$ and $N$.

We express $D_6$ as the inner direct product of $M \times N$ where $M = 
\langle r^3 \rangle$ and $N = \langle r^2, s^1\rangle$. The intersection
of $M$ and $N$ is trivial and both are normal subgroups in $D_6$.

\noindent(b) Using part (a), express $D_6$ as isomorphic to the (external) direct
product of two familiar (named) nontrival subgroups, describes the images
of $r$ and $s$ under this isomorphism.

$D_6$ can be expressed as the external direct product of $S_3 \times \mathbb{Z}_2$.
Under this isomorphism $\varphi$ we see that Im$(\varphi(s)) = ((1\text{ }2),0)$ and
Im$(\varphi(r)) = ((1\text{ }2\text{ }3),1)$.

\end{document}